\markboth{CONCLUSION AND OUTLOOK}{CONCLUSION AND OUTLOOK}
\section{Conclusion and Outlook}\label{sec:conclusion}

\subsection{Summary}\label{sec:conclusion:summary}

This paper presented Training Copilot, an open-source, multi-agent sports analytics platform that unifies heterogeneous fitness data behind a single conversational interface, addressing data fragmentation across wearable ecosystems, the lack of cross-domain contextual reasoning, and the rigid architectures of current platforms. The MCP-based integration layer achieves single-line extensibility, verified empirically by adding the Telegram proxy and flythrough servers during development. The LangGraph-based multi-agent system coordinates five domain-scoped specialists via A2A, with parallel orchestration keeping latency manageable and scoped ToolHosts reducing tool misuse; the recording-and-clipping pattern decouples data fidelity from LLM context constraints, and a RAG-based fitness specialist grounds advice in verifiable published sources with explicit citation preservation. The automated evaluation framework combines LLM judges with a deterministic trace-based grounding scorer for reproducible quality assessment.

\subsection{Future Work}\label{sec:conclusion:next}

\textbf{Field study.} The most important next step is a controlled study with real athletes: recruiting 10--20 participants from the KIT sports community, deploying Training Copilot as their training assistant for four weeks, and measuring both quantitative metrics (training consistency, load progression) and qualitative feedback (perceived usefulness, trust). The automated evaluation (Section~\ref{sec:evaluation}) assesses conversational quality; only longitudinal human evaluation can measure impact on actual training decisions.

\textbf{Knowledge base expansion.} The RAG corpus comprises fifty public-domain books, all necessarily historical; incorporating contemporary open-access publications (e.g.\ from PubMed Central) would improve relevance and let the fitness specialist cite current research. Training Copilot currently gives session-level recommendations but does not generate multi-week periodised plans~\citep{grgic_periodization_2017}; integrating plan generation with load monitoring and calendar writes is a natural extension.

\textbf{Multi-tenant deployment.} Production deployment requires per-user credential vaults (replacing the current plaintext \texttt{.tokens/} directory), session-scoped ToolHost instances, and the security measures outlined in Section~\ref{sec:discussion:limitations}. The stateless architecture (no server holds per-request state) supports horizontal scaling; the orchestrator's in-process trace assembly is the main bottleneck for this.

\textbf{Ecosystem evolution.} Strava's announcement of an official MCP server~\citep{strava_api_update_2026} validates the architecture's bet on protocol-based integration and suggests first-party MCP servers from data providers will progressively replace third-party API wrappers. As MCP adoption grows~\citep{anthropic_mcp_announcement_2024}, community-contributed servers (nutrition, air quality, social rides) could be added with a single URL; enabling token-level SSE streaming and caching frequently repeated tool calls would further reduce the 8--15 second latency discussed in Section~\ref{sec:discussion:limitations}.

\textbf{Experience-aware route planning.} Route generation already produces loops and point-to-point routes of any target distance over open map data, even where no pre-planned route exists (Section~\ref{sec:data:other}); what it does not yet do is enrich a route with points of interest along the way. Integrating a places API such as Google Places would let the route specialist answer experiential requests like ``plan a scenic bike loop where I can swim in a lake around halfway and stop for ice cream on the way back''---selecting and sequencing waypoints (caf\'es, lakes, viewpoints, restaurants) rather than optimising distance and elevation alone. Combined with a higher-quality route source (Section~\ref{sec:discussion:limitations}), this would match consumer route planners on experience while keeping the multi-source agentic reasoning that sets Training Copilot apart.

\subsection{Closing Remarks}\label{sec:conclusion:closing}

Training Copilot demonstrates that protocol-based data integration (MCP), multi-agent coordination (A2A), and domain-scoped specialist agents (LangGraph ReAct) can bridge the gap between fragmented fitness data and actionable training insight. The system runs in real time on commodity hardware and is designed to be extended by anyone who can write a Python function and a one-line configuration entry. As standardised agent protocols mature and wearable ecosystems become more interoperable, architectures like Training Copilot's---where intelligence sits in the coordination layer, not the data layer---offer a viable pattern for agentic applications beyond sports analytics.
